\documentclass[12pt,a4paper]{article}
\usepackage[utf8]{inputenc}
\usepackage{amsmath}
\usepackage{amssymb}
\usepackage{graphicx}
\usepackage{hyperref}
\usepackage{listings}
\usepackage{xcolor}
\usepackage{geometry}
\geometry{margin=1in}

\title{Differential Privacy Parameters in Federated Learning DDoS Detection}
\author{Federated Learning Project}
\date{\today}

\begin{document}

\maketitle

\section{Introduction}

This document explains the Differential Privacy (DP) parameters used in our Federated Learning system for DDoS detection. We implement DP-SGD (Differentially Private Stochastic Gradient Descent) to protect individual training samples while maintaining model utility.

\section{Differential Privacy Parameters}

\subsection{Overview}

Our implementation uses four key parameters to control the privacy-accuracy trade-off:

\begin{itemize}
    \item \textbf{Clip Norm} ($C$): Gradient clipping threshold
    \item \textbf{Noise Multiplier} ($z$): Controls noise magnitude
    \item \textbf{Delta} ($\delta$): Failure probability
    \item \textbf{Epsilon} ($\epsilon$): Privacy budget (calculated)
\end{itemize}

\subsection{Clip Norm (dp\_clip\_norm = 1.0)}

\subsubsection{Definition}

Clip Norm, denoted as $C$, is the maximum L2 norm allowed for gradient vectors. It bounds the sensitivity of gradients to individual training samples.

\subsubsection{Mathematical Formulation}

For a gradient vector $\mathbf{g}$, clipping is performed as:

\begin{equation}
\mathbf{g}_{\text{clipped}} = \mathbf{g} \cdot \min\left(1, \frac{C}{\|\mathbf{g}\|_2}\right)
\end{equation}

where $\|\mathbf{g}\|_2$ is the L2 norm of the gradient.

\subsubsection{Implementation in Our Project}

In \texttt{flower\_worker.py}, clip norm is set during worker initialization:

\begin{lstlisting}[language=Python, basicstyle=\small]
dp_clip_norm: float = 1.0  # Default value
\end{lstlisting}

The DP-SGD optimizer uses this value:

\begin{lstlisting}[language=Python, basicstyle=\small]
optimizer = DPKerasAdamOptimizer(
    l2_norm_clip=self.dp_clip_norm,  # C = 1.0
    noise_multiplier=self.dp_noise_multiplier,
    num_microbatches=1,
    learning_rate=0.001
)
\end{lstlisting}

\subsubsection{Impact}

\begin{itemize}
    \item \textbf{Lower $C$ (e.g., 0.5)}: Stronger privacy, slower convergence, lower accuracy
    \item \textbf{Higher $C$ (e.g., 2.0)}: Weaker privacy, faster convergence, higher accuracy
    \item \textbf{Our setting ($C = 1.0$)}: Balanced trade-off
\end{itemize}

\subsection{Noise Multiplier (dp\_noise\_multiplier = 1.0)}

\subsubsection{Definition}

Noise Multiplier, denoted as $z$, controls the magnitude of Gaussian noise added to gradients. It determines the standard deviation of the noise distribution.

\subsubsection{Mathematical Formulation}

After clipping, Gaussian noise is added:

\begin{equation}
\mathbf{g}_{\text{noisy}} = \mathbf{g}_{\text{clipped}} + \mathcal{N}(0, (C \cdot z)^2 \mathbf{I})
\end{equation}

where:
\begin{itemize}
    \item $C$ is the clip norm
    \item $z$ is the noise multiplier
    \item $\mathcal{N}(0, \sigma^2)$ is a Gaussian distribution with mean 0 and variance $\sigma^2$
\end{itemize}

With our settings ($C = 1.0$, $z = 1.0$), noise variance is:

\begin{equation}
\sigma^2 = (1.0 \times 1.0)^2 = 1.0
\end{equation}

\subsubsection{Implementation}

\begin{lstlisting}[language=Python, basicstyle=\small]
dp_noise_multiplier: float = 1.0  # Default value

optimizer = DPKerasAdamOptimizer(
    l2_norm_clip=self.dp_clip_norm,
    noise_multiplier=self.dp_noise_multiplier,  # z = 1.0
    ...
)
\end{lstlisting}

\subsubsection{Impact}

\begin{itemize}
    \item \textbf{Lower $z$ (e.g., 0.5)}: Less noise, weaker privacy, higher accuracy
    \item \textbf{Higher $z$ (e.g., 2.0)}: More noise, stronger privacy, lower accuracy
    \item \textbf{Our setting ($z = 1.0$)}: Standard noise level
\end{itemize}

\subsection{Delta (dp\_delta = 1e-5)}

\subsubsection{Definition}

Delta ($\delta$) is the failure probability in the $(\epsilon, \delta)$-differential privacy guarantee. It represents the probability that the privacy guarantee fails.

\subsubsection{Mathematical Meaning}

The $(\epsilon, \delta)$-DP guarantee states:

\begin{equation}
\mathbb{P}[M(D) \in S] \leq e^{\epsilon} \cdot \mathbb{P}[M(D') \in S] + \delta
\end{equation}

where:
\begin{itemize}
    \item $M$ is the mechanism (our training algorithm)
    \item $D$ and $D'$ are datasets differing by one sample
    \item $S$ is any subset of possible outputs
    \item $\delta$ is the failure probability
\end{itemize}

\subsubsection{Implementation}

\begin{lstlisting}[language=Python, basicstyle=\small]
dp_delta: float = 1e-5  # Default: 0.00001 = 0.001%

epsilon_this_round, _ = compute_dp_sgd_privacy(
    n=num_samples,
    batch_size=self.batch_size,
    noise_multiplier=self.dp_noise_multiplier,
    epochs=self.epochs_per_round,
    delta=self.dp_delta  # δ = 1e-5
)
\end{lstlisting}

\subsubsection{Interpretation}

With $\delta = 10^{-5}$:
\begin{itemize}
    \item Privacy guarantee holds with probability $1 - \delta = 99.999\%$
    \item Failure probability: 0.001\%
    \item Standard choice for datasets with $\sim 10^5$ samples
\end{itemize}

\subsubsection{Common Values}

\begin{itemize}
    \item $\delta = 1/n$: Where $n$ is the number of training samples
    \item $\delta = 10^{-5}$: Standard for medium datasets (our choice)
    \item $\delta = 10^{-6}$: Stricter, for larger datasets
\end{itemize}

\subsection{Epsilon (ε) - Privacy Budget}

\subsubsection{Definition}

Epsilon ($\epsilon$) is the privacy loss parameter. It quantifies how much information about individual samples can be inferred from the model output.

\subsubsection{Calculation}

Epsilon is calculated using the DP-SGD privacy accounting formula:

\begin{equation}
\epsilon = f(n, \text{batch\_size}, z, \text{epochs}, \delta)
\end{equation}

where:
\begin{itemize}
    \item $n$: Number of training samples
    \item $\text{batch\_size}$: Samples per batch (32 in our case)
    \item $z$: Noise multiplier (1.0)
    \item $\text{epochs}$: Training epochs per round (5)
    \item $\delta$: Failure probability (1e-5)
\end{itemize}

\subsubsection{Implementation}

In our code, epsilon is computed after each training round:

\begin{lstlisting}[language=Python, basicstyle=\small]
epsilon_this_round, _ = compute_dp_sgd_privacy(
    n=num_samples,              # e.g., 41,738 samples
    batch_size=self.batch_size,  # 32
    noise_multiplier=self.dp_noise_multiplier,  # 1.0
    epochs=self.epochs_per_round,  # 5
    delta=self.dp_delta  # 1e-5
)
self.epsilon_spent += epsilon_this_round
\end{lstlisting}

\subsubsection{Current Values}

With our settings:
\begin{itemize}
    \item Samples per worker: $\sim 41,738$
    \item Batch size: 32
    \item Epochs per round: 5
    \item Noise multiplier: 1.0
    \item \textbf{Epsilon per round}: $\approx 0.7$
    \item \textbf{Total epsilon (5 rounds)}: $\approx 3.5$
\end{itemize}

\subsubsection{Interpretation}

\begin{itemize}
    \item $\epsilon = 0$: Perfect privacy (no information released)
    \item $\epsilon = 1$: Very strong privacy
    \item $\epsilon = 3-5$: Moderate privacy (our range)
    \item $\epsilon = 10$: Weak privacy
    \item $\epsilon > 100$: Essentially no privacy
\end{itemize}

\section{How Parameters Work Together}

\subsection{DP-SGD Algorithm Flow}

The complete DP-SGD process in our implementation:

\begin{enumerate}
    \item \textbf{Forward Pass}: Compute predictions and loss
    \item \textbf{Backward Pass}: Compute gradients $\mathbf{g}$
    \item \textbf{Clipping}: $\mathbf{g}_{\text{clipped}} = \text{clip}(\mathbf{g}, C)$
    \item \textbf{Noise Addition}: $\mathbf{g}_{\text{noisy}} = \mathbf{g}_{\text{clipped}} + \mathcal{N}(0, (C \cdot z)^2)$
    \item \textbf{Update}: $\mathbf{w}_{t+1} = \mathbf{w}_t - \eta \cdot \mathbf{g}_{\text{noisy}}$
    \item \textbf{Privacy Accounting}: Compute $\epsilon$ spent
\end{enumerate}

\subsection{Privacy-Accuracy Trade-off}

The relationship between parameters and outcomes:

\begin{table}[h]
\centering
\begin{tabular}{|l|c|c|c|}
\hline
\textbf{Parameter Change} & \textbf{Privacy} & \textbf{Accuracy} & \textbf{Convergence} \\
\hline
Increase Clip Norm ($C \uparrow$) & $\downarrow$ & $\uparrow$ & $\uparrow$ \\
Decrease Clip Norm ($C \downarrow$) & $\uparrow$ & $\downarrow$ & $\downarrow$ \\
Increase Noise Mult ($z \uparrow$) & $\uparrow$ & $\downarrow$ & $\downarrow$ \\
Decrease Noise Mult ($z \downarrow$) & $\downarrow$ & $\uparrow$ & $\uparrow$ \\
\hline
\end{tabular}
\caption{Parameter Impact on System Performance}
\end{table}

\section{Implementation in Our Project}

\subsection{Worker-Side DP}

DP is implemented at the worker level in \texttt{flower\_worker.py}:

\begin{lstlisting}[language=Python, basicstyle=\small]
class FlowerWorker(fl.client.NumPyClient):
    def __init__(self, ..., enable_dp=False, 
                 dp_clip_norm=1.0, 
                 dp_noise_multiplier=1.0, 
                 dp_delta=1e-5):
        self.enable_dp = enable_dp and DP_AVAILABLE
        self.dp_clip_norm = dp_clip_norm
        self.dp_noise_multiplier = dp_noise_multiplier
        self.dp_delta = dp_delta
        self.epsilon_spent = 0.0
\end{lstlisting}

\subsection{DP Optimizer Usage}

When DP is enabled, the model uses DP-SGD optimizer:

\begin{lstlisting}[language=Python, basicstyle=\small]
if self.enable_dp:
    optimizer = DPKerasAdamOptimizer(
        l2_norm_clip=self.dp_clip_norm,      # C = 1.0
        noise_multiplier=self.dp_noise_multiplier,  # z = 1.0
        num_microbatches=1,
        learning_rate=0.001
    )
else:
    optimizer = Adam(learning_rate=0.001)
\end{lstlisting}

\subsection{Privacy Budget Tracking}

After each training round, epsilon is computed and accumulated:

\begin{lstlisting}[language=Python, basicstyle=\small]
if self.enable_dp and compute_dp_sgd_privacy is not None:
    epsilon_this_round, _ = compute_dp_sgd_privacy(
        n=num_samples,
        batch_size=self.batch_size,
        noise_multiplier=self.dp_noise_multiplier,
        epochs=self.epochs_per_round,
        delta=self.dp_delta
    )
    self.epsilon_spent += epsilon_this_round
\end{lstlisting}

\subsection{Environment Variable Control}

DP can be enabled/disabled via environment variable:

\begin{lstlisting}[language=bash, basicstyle=\small]
# Enable DP
FL_ENABLE_DP=true docker compose up -d

# Disable DP (default)
FL_ENABLE_DP=false docker compose up -d
\end{lstlisting}

\section{Privacy Budget Accumulation}

\subsection{Per-Round Epsilon}

Each training round consumes privacy budget:

\begin{equation}
\epsilon_{\text{round}} = f(n, B, z, E, \delta)
\end{equation}

where:
\begin{itemize}
    \item $n$: Samples per worker (41,738)
    \item $B$: Batch size (32)
    \item $z$: Noise multiplier (1.0)
    \item $E$: Epochs per round (5)
    \item $\delta$: Delta (1e-5)
\end{itemize}

\subsection{Total Privacy Budget}

After $R$ rounds:

\begin{equation}
\epsilon_{\text{total}} = \sum_{r=1}^{R} \epsilon_{\text{round}} = R \times \epsilon_{\text{round}}
\end{equation}

For our 5-round training:
\begin{equation}
\epsilon_{\text{total}} \approx 5 \times 0.7 = 3.5
\end{equation}

\section{Example: Privacy Budget Calculation}

\subsection{Scenario}

\begin{itemize}
    \item Worker 1: 41,738 samples
    \item Batch size: 32
    \item Epochs per round: 5
    \item Noise multiplier: 1.0
    \item Delta: 1e-5
    \item Total rounds: 5
\end{itemize}

\subsection{Calculation}

Using the DP-SGD privacy accounting:

\begin{align}
\epsilon_{\text{round}} &\approx 0.6979 \\
\epsilon_{\text{total}} &= 5 \times 0.6979 = 3.4895
\end{align}

\subsection{Interpretation}

After 5 rounds of federated training:
\begin{itemize}
    \item Privacy guarantee: $(3.49, 10^{-5})$-DP
    \item Meaning: Adding/removing one sample changes output by factor $e^{3.49} \approx 32.8$
    \item Failure probability: 0.001\%
\end{itemize}

\section{Recommendations}

\subsection{For Stronger Privacy}

\begin{itemize}
    \item Reduce clip norm: $C = 0.5$
    \item Increase noise multiplier: $z = 2.0$
    \item Reduce epochs per round: $E = 3$
    \item Result: $\epsilon \approx 0.3$ per round, total $\approx 1.5$
\end{itemize}

\subsection{For Better Accuracy}

\begin{itemize}
    \item Increase clip norm: $C = 2.0$
    \item Reduce noise multiplier: $z = 0.5$
    \item Keep epochs: $E = 5$
    \item Result: $\epsilon \approx 1.2$ per round, total $\approx 6.0$
\end{itemize}

\subsection{For Balanced Trade-off}

\begin{itemize}
    \item Current settings: $C = 1.0$, $z = 1.0$, $E = 5$
    \item Result: $\epsilon \approx 0.7$ per round, total $\approx 3.5$
    \item Good balance between privacy and accuracy
\end{itemize}

\section{Conclusion}

Our DP implementation uses four key parameters to control the privacy-accuracy trade-off:

\begin{itemize}
    \item \textbf{Clip Norm} ($C = 1.0$): Bounds gradient sensitivity
    \item \textbf{Noise Multiplier} ($z = 1.0$): Controls noise magnitude
    \item \textbf{Delta} ($\delta = 10^{-5}$): Failure probability
    \item \textbf{Epsilon} ($\epsilon \approx 0.7$/round): Privacy budget
\end{itemize}

These parameters work together to provide $(3.5, 10^{-5})$-differential privacy after 5 rounds of federated training, protecting individual training samples while maintaining reasonable model accuracy.

\end{document}

